% page 352 of the facsimile (image p-351 of the scan)
% block 167.6 x 246.3 mm ; left margin 8.0 mm
\pgc{-12.700mm}{0.000mm}{
\l{\cel{4}344.}
\l{}
\l{}
\l{\cel{1}lokomocar. (netrans.)\cel{2}Agar la movo o la serio de movi per}
\l{\cel{4}qui onu su transportas de loko ad altra loko.- DEFIS.}
\l{}
\l{\cel{1}lokomotivo. I, (olim).Vaporo-mashino portata per roti, ed}
\l{\cel{4}uzata por tirar treni sur relvoyi. II. (\sou{nun}) Vehilo}
\l{\cel{4}analoga, ma movata da elektro o da Diesel-motoro.}
\l{\cel{4}- DEFIRS.}
\l{}
\l{\cel{1}\sou{lokusto}. (\sou{zool.}) Insekto aloza, ek la familio \textquotedbl{}ortopteri\textquotedbl{},}
\l{\cel{4}saltera. - L.\cel{1}\sou{locusta}. - EFIS.}
\l{}
\l{\cel{1}lolio. (bot.)\cel{2}Herbo di qua la grani esas nigra, qua kreskas}
\l{\cel{4}sur frumento. - L. lolium temulentum.- EIL.}
\l{}
\l{\cel{1}\sou{lombardo}. Establisuro di presto po gajo (institucita kun}
\l{\cel{4}skopo karitatala). - DEF.}
\l{}
\l{\cel{1}\sou{lombriko}. (\sou{zool.)}\cel{2}Anelideo, di qua la nomo esas anke ter-}
\l{\cel{4}vermo. L.\cel{1}\sou{lumbricus terrestris}. - EFIS.}
\l{}
\l{\cel{1}\sou{loncho.}\cel{2}Porciono plu o min dina di korpo, tranchita per}
\l{\cel{4}de-seko tre egala, tale ke la surfaco esas tre plana.}
\l{}
\l{}
\l{\cel{1}\sou{lon}ga. I. (\sou{spaco)}. Qua esas extensita, de un de sua extre-}
\l{\cel{4}majo a la altra, plumulte pri ta dimensiono kam pri la}
\l{\cel{4}cetera. - Qua esas (tante e tante) extensita pri ta dimen-}
\l{\cel{4}siono. (ex. Ta domo esas longa ye X metri). -II. (tempo).}
\l{\cel{4}Qua esas tre extensita en tempo. - DEFI.}
\l{}
\l{\sou{longitudo}. I. (\sou{geogr.}) Koordinato geografiala qua uzesas por}
\l{\cel{4}determinar la situeso di loko sur la surfaco di la Ter-}
\l{\cel{4}globo : disto de ta loko a la meridiano determinita.}
\l{\cel{4}- II. (\sou{astron.}) Koordinato astronomiala uzata por determi-}
\l{\cel{4}nar la situeso di astro sur la sfero cielala ; disto de}
\l{\cel{4}ta astro a la punto equinoxala di printempo. - EFIRS.}
\l{}
\l{\cel{1}\sou{lor}. En la epoko kande... En la epoko di... (\sou{lore...lore}\cel{1}:}
\l{\cel{4}ye ula instanto... ye altra...).}
\l{}
\l{\cel{1}\sou{lorno}. Aparato qua konsistas ye lensi quin onu interpozas in-}
\l{\cel{4}ter la okulo regardanta e la objekto regardata, por vidar}
\l{\cel{4}plu precize, plu dicerneble. - FR.}
\l{}
\l{\cel{1}\sou{loto}. To quo atribuesas ad ulu lor partigo o disdono. To quo}
\l{\cel{4}atribuesas da hazardo ad ulu lor lotrio. - EF.}
\l{}
\l{\cel{1}\sou{lotao}.\cel{2}(\sou{zool.}) Fisho qua vivas en aquo sen-sala, ek la gene-}
\l{\cel{4}ro \textquotedbl{}gado\textquotedbl{}. - L. gadus lota. - EFSL.}
\l{}
\l{\cel{1}\sou{lotriar}. (\sou{trans.}) Agar ta hazardo-ludo en qua onu distributas}
\l{\cel{4}ula quanto de bilieti pagenda, qui havas numeri, ek qui uli,}
\l{\cel{4}indikota da hazardo (fato), yurizas pri recevor pekunio-quanto}
\l{\cel{4}od objekto. - DEFIRS.}
}
